\pagestyle{empty}
\tight
\begin{tblr}{width=\textwidth,colspec={|Q[c,m,1.9cm]|X[l,m]|},hlines,vlines,hspan=minimal,row{1}={bg=headblue},rowsep=1pt,colsep=3pt}
\SetCell{c,m}\hfil\logo\hfil & \textbf{POLITEKNIK NEGERI MALANG}\par\textbf{JURUSAN TEKNOLOGI INFORMASI}\par\textbf{PROGRAM STUDI: D4 TEKNIK INFORMATIKA} \\
\SetCell[c=2]{c} \textbf{RENCANA PEMBELAJARAN SEMESTER (RPS)} & \\
\end{tblr}
\begin{longtblr}{width=\textwidth,colspec={|Q[l,m,1.9cm]|X[0.85,l,m]|X[1.45,l,m]|X[0.75,l,m]|X[1.15,l,m]|},hlines,vlines,hspan=minimal,rowsep=1pt,colsep=2pt,row{1,3}={bg=headgray,font=\bfseries}}
MATA KULIAH & KODE & BOBOT (SKS)/jam & SEMESTER & TGL. PENYUSUNAN \\
Internet of Things & RTI256105 & 3 SKS/6 jam & 6 & 1 Juli 2025 \\
OTORISASI & Dosen Pengembang RPS & Koordinator RMK & \SetCell[c=2]{l} Ka PRODI &  \\
 & Vipkas Al Hadid Firdaus, S.T., M.T. &  & \SetCell[c=2]{l}{\vspace{8mm}\par Dr. Ely Setyo Astuti, S.T., M.T.} &  \\
\SetCell[r=13]{l} Capaian Pembelajaran (CP) & \SetCell[c=4]{l,bg=headgray,font=\bfseries} Capaian Pembelajaran Lulusan Program Studi (CPL-Prodi) &  &  &  \\
 & CPL05 & \SetCell[c=3]{l} Mampu menghasilkan solusi teknologi komputasi dan infrastruktur digital yang sesuai dengan kebutuhan organisasi dan pengguna. &  &  \\
 & CPL02 & \SetCell[c=3]{l} Mampu menerapkan prinsip rekayasa sistem dan perangkat lunak untuk menghasilkan solusi aplikasi multi-platform sesuai kebutuhan. &  &  \\
 & CPL09 & \SetCell[c=3]{l} Mampu menjelaskan cara kerja sistem komputer dan menerapkan pendekatan komputasi untuk menyelesaikan masalah. &  &  \\
 & \SetCell[c=4]{l,bg=headgray,font=\bfseries} Capaian Pembelajaran mata kuliah (CPL-MK) &  &  &  \\
 & CPMK0204 & \SetCell[c=3]{l} Mahasiswa mampu menganalisis kebutuhan komputasi dan merancang arsitektur solusi IoT terdistribusi. &  &  \\
 & CPMK0503 & \SetCell[c=3]{l} Mahasiswa mampu mengimplementasikan solusi komputasi paralel-terdistribusi untuk sistem IoT. &  &  \\
 & CPMK0904 & \SetCell[c=3]{l} Mahasiswa mampu menguraikan prinsip kerja sistem paralel-terdistribusi yang dimanfaatkan untuk pemrosesan data skala besar. &  &  \\
 & \SetCell[c=4]{l,bg=headgray,font=\bfseries} Kemampuan akhir tiap tahapan belajar (Sub-CPMK) &  &  &  \\
 & SCPMK0204-04901 & \SetCell[c=3]{l} Mahasiswa mampu merancang arsitektur sistem IoT end-to-end meliputi sensor, gateway, dan cloud platform. &  &  \\
 & SCPMK0503-04902 & \SetCell[c=3]{l} Mahasiswa mampu mengimplementasikan protokol komunikasi IoT (MQTT, HTTP, CoAP) dan pemrosesan data sensor terdistribusi. &  &  \\
 & SCPMK0904-04903 & \SetCell[c=3]{l} Mahasiswa mampu mengkonfigurasi platform IoT (Node-RED, AWS IoT, atau setara) dan mengintegrasikan layanan cloud untuk monitoring dan kontrol perangkat. &  &  \\
\SetCell[r=2]{l} Penilaian & \SetCell[c=4]{l} *Nilai CPMK didapat dari agregasi nilai SUB-CPMK yang dibebankan &  &  &  \\
 & \SetCell[c=4]{l}{\tiny\begin{tblr}{width=\linewidth,colspec={|Q[c,m,0.1\linewidth]|Q[c,m,0.16\linewidth]|X[l,m]|X[l,m]|X[l,m]|X[l,m]|X[l,m]|X[l,m]|Q[c,m,0.08\linewidth]|},hlines,vlines,hspan=minimal,rowsep=1pt,colsep=0.7pt,row{1,2}={bg=headgray,font=\bfseries}}
Kode CPMK & Kode SCPMK & \SetCell[c=6]{c} Bobot per Bentuk Penilaian & & & & & & Total Bobot \\
 & & Tugas & Kuis 1 & UTS & Kuis 2 & PBL & UAS & \\
CPMK0204 & SCPMK0204-04901 & 14\%: perancangan arsitektur IoT & 5\%: konsep dan arsitektur IoT & 8\%: arsitektur sistem IoT & 0\% & 0\% & 0\% & 27\% \\
CPMK0503 & SCPMK0503-04902 & 14\%: implementasi protokol IoT & 0\% & 7\%: protokol dan komunikasi IoT & 5\%: platform dan protokol IoT & 12\%: implementasi sistem IoT & 0\% & 38\% \\
CPMK0904 & SCPMK0904-04903 & 0\% & 0\% & 0\% & 0\% & 13\%: integrasi cloud dan monitoring & 22\%: demo sistem IoT & 35\% \\
\SetCell[c=8]{r}\textbf{Total Per Penilaian} & & & & & & & & \textbf{100\%} \\
\end{tblr}} &  &  &  \\
Diskripsi Singkat Mata Kuliah & \SetCell[c=4]{l} Internet of Things membangun kompetensi merancang, mengimplementasikan, dan mendeploy sistem IoT end-to-end dari sensor hingga cloud menggunakan protokol standar industri dan platform IoT modern. &  &  &  \\
Bahan Kajian / Materi Pembelajaran & \SetCell[c=4]{l}\begin{enumerate}\item Arsitektur IoT: konsep, komponen, sensor, aktuator, gateway, dan layer edge-fog-cloud\item Mikrokontroler dan SBC: Arduino, Raspberry Pi, dan antarmuka sensor\item Protokol komunikasi IoT: MQTT, HTTP, CoAP, WebSocket, dan pemrosesan data sensor\item Platform IoT: Node-RED, AWS IoT, keamanan, integrasi cloud, dan visualisasi data real-time\end{enumerate} &  &  &  \\
\SetCell[r=2]{l} Pustaka & \SetCell[c=4]{l} \textbf{Utama:}\begin{enumerate}\item Buyya, R. \& Dastjerdi, A. 2016. Internet of Things: Principles and Paradigms. Elsevier.\item McEwen, A. 2014. Designing the Internet of Things. Wiley.\item Dokumentasi AWS IoT Core dan Node-RED.\end{enumerate} &  &  &  \\
 & \SetCell[c=4]{l} \textbf{Pendukung:} Materi LMS, dokumentasi resmi platform IoT, dan jobsheet praktikum. &  &  &  \\
Media Pembelajaran & \SetCell[c=2]{l} \textbf{Software:} Arduino IDE, Node-RED, MQTT broker (Mosquitto), AWS IoT Console, Fritzing, LMS. &  & \SetCell[c=2]{l} \textbf{Hardware:} Komputer/laptop, mikrokontroler (Arduino/Raspberry Pi), sensor kit, proyektor. &  \\
Nama Dosen Pengampu & \SetCell[c=4]{l} Tim Pengajar Program Studi D4 TI &  &  &  \\
Matakuliah Syarat & \SetCell[c=4]{l} - &  &  &  \\
\end{longtblr}
\begin{longtblr}{width=\textwidth,colspec={|Q[c,m,0.06\textwidth]|X[1.35,l,m]|X[1.08,l,m]|X[1.16,l,m]|X[0.7,l,m]|X[1.24,l,m]|X[1.06,l,m]|X[0.98,l,m]|Q[c,m,0.05\textwidth]|},hlines,vlines,hspan=minimal,rowhead=2,row{1,2}={bg=headgreen,font=\bfseries},rowsep=1pt,colsep=1.5pt}
Mgg.\par Ke & Kemampuan Akhir Yang Direncanakan (Sub-CP-MK) & Bahan kajian (Materi Pembelajaran) & Bentuk dan Metode Pembelajaran & Estimasi Waktu & Pengalaman Belajar Mahasiswa & Kriteria \& Bentuk Penilaian & Indikator Penilaian & Bobot Penilaian (\%) \\
(1) & (2) & (3) & (4) & (5) & (6) & (7) & (8) & (9) \\
1 & SCPMK0204-04901 & Pengantar IoT: konsep, sejarah, dan ekosistem. & Modalitas: Blended Learning. Bentuk: Luring/Daring. Metode: Case Method, praktik IoT, studi kasus, dan peer review. & 1 x 3 x 50' tatap muka; 1 x 3 x 50' tugas/praktik mandiri. & Mahasiswa mengerjakan aktivitas terarah untuk pengantar IoT: konsep, sejarah, dan ekosistem dan menunjukkan evidence hasil belajar. & Bentuk penilaian: Pengantar IoT. Mengacu rubrik RTM. & Ketepatan konsep; kualitas implementasi; validasi dan dokumentasi. & 2 \\
2 & SCPMK0204-04901 & Arsitektur IoT: edge, fog, dan cloud layers. & Modalitas: Blended Learning. Bentuk: Luring/Daring. Metode: Case Method, praktik IoT, studi kasus, dan peer review. & 1 x 3 x 50' tatap muka; 1 x 3 x 50' tugas/praktik mandiri. & Mahasiswa mengerjakan aktivitas terarah untuk arsitektur IoT dan menunjukkan evidence hasil belajar. & Bentuk penilaian: Arsitektur IoT. Mengacu rubrik RTM. & Ketepatan konsep; kualitas implementasi; validasi dan dokumentasi. & 3 \\
3 & SCPMK0204-04901 & Sensor dan aktuator: karakteristik dan antarmuka. & Modalitas: Blended Learning. Bentuk: Luring/Daring. Metode: Case Method, praktik IoT, studi kasus, dan peer review. & 1 x 3 x 50' tatap muka; 1 x 3 x 50' tugas/praktik mandiri. & Mahasiswa mengerjakan aktivitas terarah untuk sensor dan aktuator dan menunjukkan evidence hasil belajar. & Bentuk penilaian: Sensor dan Aktuator. Mengacu rubrik RTM. & Ketepatan konsep; kualitas implementasi; validasi dan dokumentasi. & 3 \\
4 & SCPMK0204-04901 & Kuis 1 arsitektur dan komponen IoT. & Modalitas: Blended Learning. Bentuk: Luring/Daring. Metode: Case Method, praktik IoT, studi kasus, dan peer review. & 1 x 3 x 50' tatap muka; 1 x 3 x 50' tugas/praktik mandiri. & Mahasiswa mengerjakan aktivitas terarah untuk kuis 1 arsitektur dan komponen IoT dan menunjukkan evidence hasil belajar. & Bentuk penilaian: Kuis 1 Arsitektur IoT. Mengacu rubrik RTM. & Ketepatan konsep; kualitas implementasi; validasi dan dokumentasi. & 5 \\
5 & SCPMK0503-04902 & Mikrokontroler dan SBC: Arduino, Raspberry Pi. & Modalitas: Blended Learning. Bentuk: Luring/Daring. Metode: Case Method, praktik IoT, studi kasus, dan peer review. & 1 x 3 x 50' tatap muka; 1 x 3 x 50' tugas/praktik mandiri. & Mahasiswa mengerjakan aktivitas terarah untuk mikrokontroler dan SBC dan menunjukkan evidence hasil belajar. & Bentuk penilaian: Mikrokontroler dan SBC. Mengacu rubrik RTM. & Ketepatan konsep; kualitas implementasi; validasi dan dokumentasi. & 2 \\
6 & SCPMK0503-04902 & Protokol komunikasi: MQTT, HTTP, CoAP, WebSocket. & Modalitas: Blended Learning. Bentuk: Luring/Daring. Metode: Case Method, praktik IoT, studi kasus, dan peer review. & 1 x 3 x 50' tatap muka; 1 x 3 x 50' tugas/praktik mandiri. & Mahasiswa mengerjakan aktivitas terarah untuk protokol komunikasi IoT dan menunjukkan evidence hasil belajar. & Bentuk penilaian: Protokol Komunikasi IoT. Mengacu rubrik RTM. & Ketepatan konsep; kualitas implementasi; validasi dan dokumentasi. & 2 \\
7 & SCPMK0503-04902 & Gateway IoT dan edge computing. & Modalitas: Blended Learning. Bentuk: Luring/Daring. Metode: Case Method, praktik IoT, studi kasus, dan peer review. & 1 x 3 x 50' tatap muka; 1 x 3 x 50' tugas/praktik mandiri. & Mahasiswa mengerjakan aktivitas terarah untuk gateway IoT dan edge computing dan menunjukkan evidence hasil belajar. & Bentuk penilaian: Gateway IoT dan Edge Computing. Mengacu rubrik RTM. & Ketepatan konsep; kualitas implementasi; validasi dan dokumentasi. & 2 \\
8 & SCPMK0204-04901, SCPMK0503-04902 & UTS perancangan arsitektur sistem IoT. & Modalitas: Blended Learning. Bentuk: Luring/Daring. Metode: Case Method, praktik IoT, studi kasus, dan peer review. & 1 x 3 x 50' tatap muka; 1 x 3 x 50' tugas/praktik mandiri. & Mahasiswa mengerjakan aktivitas terarah untuk UTS perancangan arsitektur sistem IoT dan menunjukkan evidence hasil belajar. & Bentuk penilaian: UTS Perancangan Arsitektur IoT. Mengacu rubrik RTM. & Ketepatan konsep; kualitas implementasi; validasi dan dokumentasi. & 15 \\
9 & SCPMK0904-04903 & Platform IoT: Node-RED, ThingSpeak, AWS IoT. & Modalitas: Blended Learning. Bentuk: Luring/Daring. Metode: Case Method, praktik IoT, studi kasus, dan peer review. & 1 x 3 x 50' tatap muka; 1 x 3 x 50' tugas/praktik mandiri. & Mahasiswa mengerjakan aktivitas terarah untuk platform IoT dan menunjukkan evidence hasil belajar. & Bentuk penilaian: Platform IoT. Mengacu rubrik RTM. & Ketepatan konsep; kualitas implementasi; validasi dan dokumentasi. & 3 \\
10 & SCPMK0503-04902 & Pengolahan data sensor dan time-series database. & Modalitas: Blended Learning. Bentuk: Luring/Daring. Metode: Case Method, praktik IoT, studi kasus, dan peer review. & 1 x 3 x 50' tatap muka; 1 x 3 x 50' tugas/praktik mandiri. & Mahasiswa mengerjakan aktivitas terarah untuk pengolahan data sensor dan menunjukkan evidence hasil belajar. & Bentuk penilaian: Pengolahan Data Sensor. Mengacu rubrik RTM. & Ketepatan konsep; kualitas implementasi; validasi dan dokumentasi. & 3 \\
11 & SCPMK0503-04902 & Kuis 2 protokol dan platform IoT. & Modalitas: Blended Learning. Bentuk: Luring/Daring. Metode: Case Method, praktik IoT, studi kasus, dan peer review. & 1 x 3 x 50' tatap muka; 1 x 3 x 50' tugas/praktik mandiri. & Mahasiswa mengerjakan aktivitas terarah untuk kuis 2 protokol dan platform IoT dan menunjukkan evidence hasil belajar. & Bentuk penilaian: Kuis 2 Protokol dan Platform IoT. Mengacu rubrik RTM. & Ketepatan konsep; kualitas implementasi; validasi dan dokumentasi. & 5 \\
12 & SCPMK0904-04903 & Keamanan sistem IoT. & Modalitas: Blended Learning. Bentuk: Luring/Daring. Metode: Case Method, praktik IoT, studi kasus, dan peer review. & 1 x 3 x 50' tatap muka; 1 x 3 x 50' tugas/praktik mandiri. & Mahasiswa mengerjakan aktivitas terarah untuk keamanan sistem IoT dan menunjukkan evidence hasil belajar. & Bentuk penilaian: Keamanan Sistem IoT. Mengacu rubrik RTM. & Ketepatan konsep; kualitas implementasi; validasi dan dokumentasi. & 2 \\
13 & SCPMK0904-04903 & Integrasi cloud dan visualisasi data real-time. & Modalitas: Blended Learning. Bentuk: Luring/Daring. Metode: Case Method, praktik IoT, studi kasus, dan peer review. & 1 x 3 x 50' tatap muka; 1 x 3 x 50' tugas/praktik mandiri. & Mahasiswa mengerjakan aktivitas terarah untuk integrasi cloud dan visualisasi data dan menunjukkan evidence hasil belajar. & Bentuk penilaian: Integrasi Cloud dan Visualisasi. Mengacu rubrik RTM. & Ketepatan konsep; kualitas implementasi; validasi dan dokumentasi. & 2 \\
14 & SCPMK0904-04903 & Deployment dan monitoring sistem IoT. & Modalitas: Blended Learning. Bentuk: Luring/Daring. Metode: Case Method, praktik IoT, studi kasus, dan peer review. & 1 x 3 x 50' tatap muka; 1 x 3 x 50' tugas/praktik mandiri. & Mahasiswa mengerjakan aktivitas terarah untuk deployment dan monitoring IoT dan menunjukkan evidence hasil belajar. & Bentuk penilaian: Deployment dan Monitoring IoT. Mengacu rubrik RTM. & Ketepatan konsep; kualitas implementasi; validasi dan dokumentasi. & 2 \\
15 & SCPMK0904-04903 & Studi kasus: smart home, agriculture, smart city. & Modalitas: Blended Learning. Bentuk: Luring/Daring. Metode: Case Method, praktik IoT, studi kasus, dan peer review. & 1 x 3 x 50' tatap muka; 1 x 3 x 50' tugas/praktik mandiri. & Mahasiswa mengerjakan aktivitas terarah untuk studi kasus IoT dan menunjukkan evidence hasil belajar. & Bentuk penilaian: Studi Kasus IoT. Mengacu rubrik RTM. & Ketepatan konsep; kualitas implementasi; validasi dan dokumentasi. & 2 \\
16 & SCPMK0503-04902, SCPMK0904-04903 & PBL implementasi prototipe sistem IoT. & Modalitas: Blended Learning. Bentuk: Luring/Daring. Metode: Case Method, praktik IoT, studi kasus, dan peer review. & 1 x 3 x 50' tatap muka; 1 x 3 x 50' tugas/praktik mandiri. & Mahasiswa mengerjakan aktivitas terarah untuk PBL implementasi prototipe IoT dan menunjukkan evidence hasil belajar. & Bentuk penilaian: PBL Implementasi Prototipe IoT. Mengacu rubrik RTM. & Ketepatan konsep; kualitas implementasi; validasi dan dokumentasi. & 25 \\
17 & SCPMK0904-04903 & UAS demo dan presentasi sistem IoT. & Modalitas: Blended Learning. Bentuk: Luring/Daring. Metode: Case Method, praktik IoT, studi kasus, dan peer review. & 1 x 3 x 50' tatap muka; 1 x 3 x 50' tugas/praktik mandiri. & Mahasiswa mengerjakan aktivitas terarah untuk UAS demo dan presentasi sistem IoT dan menunjukkan evidence hasil belajar. & Bentuk penilaian: UAS Demo Sistem IoT. Mengacu rubrik RTM. & Ketepatan konsep; kualitas implementasi; validasi dan dokumentasi. & 22 \\
\end{longtblr}
